\section{Extension to Multi-Layer Blocks}
\label{sec:multilayer}

\begin{figure*}[t]
  \centering
  \includegraphics[width=0.92\textwidth]{figures/lr_transfer_combined.pdf}
  \caption{\textbf{Learning-rate transfer across scaling rules and depths.}
  Curves sweep learning rate for loop counts $N\in\{1,2,4,8\}$; stars mark per-$N$ optima.
  The $x$-axis is the base learning rate $\eta_0$; the actual repeated-block learning rate is $\eta_0\,m_L^{-1/2}=\eta_0(L/12)^{-1/2}$, so for the $L{=}12$ panels (a,b) plotted and actual coincide, while panels (c,d) at $L{=}24,48$ apply the depth correction implicitly.
  At $L{=}12$, sqrt scaling shifts the optimum with $N$ (panel~a), while linear scaling aligns optima and improves large-$N$ loss (panel~b).
  With $\varepsilon=1/(N\sqrt{L})$, a single base learning rate remains near-optimal across both $L$ and $N$.
  Diverged runs (final validation loss $>4$) are omitted.}
  \label{fig:lr-transfer}
\vspace{-0.3cm}
\end{figure*}

\begin{figure*}[t]
  \centering
  \includegraphics[width=\textwidth]{figures/update_cosine.pdf}
  \caption{\textbf{Cosine similarity between loop-step increments $\delta_n=h_n-h_{n-1}$ after full training} (20{,}000 steps, $N{=}8$).
  Most early and middle adjacent loop steps remain positively correlated, and the overall pattern is qualitatively consistent across depths; pairs involving the final loop step (step~8) are weaker and can be negative.}
  \label{fig:update-cosine}
\vspace{-0.3cm}
\end{figure*}

\begin{figure*}[t]
  \centering
  \includegraphics[width=\textwidth]{figures/hidden_norm.pdf}
  \caption{\textbf{Residual-stream norm trajectories across loop steps} under linear residual scaling.
  Each line traces the residual-stream L2 norm from loop step~0 (post-embedding) through the final step, for $N\in\{2,4,8\}$.
  Although the per-step increments differ across $N$, the final norms converge to approximately the same value within each depth group, confirming that the $1/N$ rescaling keeps the residual stream bounded regardless of $N$.}
  \label{fig:hidden-norm}
\vspace{-0.3cm}
\end{figure*}


Section~\ref{sec:theory} analyzed a single shared layer; practical looped architectures repeat a block of $L$ distinct layers $N$ times.
We now extend the scaling analysis to this setting.

\subsection{Setup}

Each layer $\ell$ has its own weight matrix $W_\ell$, but the same $W_\ell$ is reused across all $N$ iterations.
We index hidden states as $h_{n,\ell}$, where $n\in\{0,\dots,N-1\}$ is the loop iteration and $\ell\in\{0,\dots,L-1\}$ is the layer index within the block.
The recursion is:
\begin{equation}
  \begin{aligned}
  h_{n,\ell+1}
  &= h_{n,\ell}+\varepsilon\, W_\ell\,\phi(h_{n,\ell}),\\
  \ell&=0,\dots,L-1,\qquad n=0,\dots,N-1,
  \end{aligned}
  \label{eq:Llayer-recursion}
\end{equation}
where the output of the last layer feeds into the next iteration: $h_{n+1,0} = h_{n,L}$, starting from $h_{0,0} = h_0$.
The model output is $h_{N-1,L}$, equivalently $h_{N,0}$ after the final loop pass; our goal is to determine $\varepsilon$ as a function of both $N$ and $L$.

\subsection{Two-source variance accumulation}

Telescoping the residual additions gives:
\[
  h_{\text{out}}
  =
  h_0+\varepsilon\sum_{\ell=0}^{L-1}G_\ell,
  \qquad
  G_\ell \triangleq W_\ell\sum_{n=0}^{N-1}\phi(h_{n,\ell}).
\]
Unlike the single-layer case, the output variance now receives contributions from two distinct sources.
Within each layer $\ell$, weight sharing produces the same quadratic accumulation as in Section~\ref{sec:theory}: successive activations $\phi(h_{n,\ell})$ are correlated across loop steps, giving $\|G_\ell\|^2=\Theta(N^2 d)$.
Across the $L$ layers, the independent weight matrices contribute a standard depth-wise random walk, the same mechanism studied for non-shared deep networks~\citep{dey2025completep}.
The only additional issue is that the $L$ layers are not literally independent: they communicate through the same residual stream.
The factorized scaling is therefore valid when this communication remains local in strength.
Informally, after the within-layer loop accumulation has been normalized by $1/N$, changing one unique layer should perturb the other normalized branch outputs only through the total residual size $\varepsilon N$, rather than dragging all layers in a single coherent direction.
Under this weak cross-layer-coupling picture, the remaining accumulation over $\ell$ is the ordinary depth-wise random walk, giving the extra $1/\!\sqrt{L}$ factor.
The post-training hidden-norm and cosine diagnostics in Figures~\ref{fig:hidden-norm} and~\ref{fig:update-cosine} are consistent with this picture: the residual stream remains controlled across loop counts, while correlations persist without producing a global layer-wise collapse.
Writing $\varepsilon=\lambda/(N\!\sqrt{L})$ with a tunable $O(1)$ constant $\lambda$, the following theorem formalizes this intuition.

\begin{theorem}[Multi-layer block scaling; informal]
\label{thm:Llayer}
Assume that\textup{:}
\begin{enumerate}
\item[\textup{(i)}] each layer's branch output has nondegenerate normalized norm\textup{;}
\item[\textup{(ii)}] the conditional mean of each layer's branch output, given the other layers' weights, has normalized norm at most $O(\varepsilon N)$\textup{;}
\item[\textup{(iii)}] replacing any single weight matrix changes other branches' outputs by at most $O(\varepsilon N)$\textup{.}
\end{enumerate}
Then $\varepsilon=\lambda/(N\!\sqrt{L})$ yields a residual-stream norm bounded by a function of $\lambda$ alone, independent of both $N$ and $L$.
For sufficiently small $\lambda$, the bound is tight\textup{:} the unscaled total branch variance is $\Theta(LN^2)$.
The formal statement and proof are in Appendix~\ref{app:Llayer}, Theorem~\ref{thm:Llayer-formal}.
\end{theorem}

\subsection{Learning-rate scaling}

Under the factorized scaling, an update $\Delta W_\ell$ to layer $\ell$ is applied at all $N$ loop steps, contributing $O(\varepsilon\eta N)=O(\eta\lambda/\!\sqrt{L})$ to the output change.
In the worst case, the $L$ layer updates add coherently with norm $O(L)$, giving total output change $O(\eta\lambda\sqrt{L})$.
Requiring $O(1)$ output change yields:
\(
  \eta \lesssim \frac{1}{\lambda\sqrt{L}}.
\)
Proposition~\ref{prop:Llayer-lr} (Appendix~\ref{app:Llayer-lr}) formalizes this as a linearized maximal-update bound, sharp under a coherence condition on the layer updates.
Crucially, $N$ does not appear: once $1/N$ scaling is applied to the residual branch, the learning rate depends only on the unique depth $L$, enabling hyperparameter transfer.
Table~\ref{tab:recipe} summarizes the resulting scaling recipe; the full per-component parameterization is in Appendix~\ref{app:completep-loop-table}.

\begin{table}[t]
\centering
\small
\begin{tabular}{@{}ll@{}}
\toprule
Component & Scaling rule \\
\midrule
Residual branch & $\times\;\lambda\,N^{-1}\,m_L^{-1/2}$ \\
Looped stack LR & $\eta_0\,m_L^{-1/2}$ \\
Embed / Unembed LR & $\eta_0$ \\
\bottomrule
\end{tabular}
\caption{
Executive scaling recipe for a looped Transformer with $L$ unique layers repeated $N$ times.
$m_L=L/L_{\mathrm{ref}}$ is the depth multiplier relative to a reference model. $\lambda$ is a tunable $O(1)$ constant.
}
\label{tab:recipe}
\vspace{-0.3cm}
\end{table}
